\documentclass[10pt,a4paper]{amsart}
\usepackage[margin=19mm]{geometry}
\usepackage{amsmath,amssymb,mathtools}
\usepackage[scr=rsfs,cal=euler]{mathalfa}
\usepackage{xfrac}
\usepackage[unicode,hidelinks,bookmarks=false]{hyperref}
\newcommand{\ie}{i.e.\ }
\newcommand{\cf}{cf.\ }
\newcommand{\resp}{respectively\ }
\newcommand{\Subsection}{\subsection}
\providecommand{\nobreakdash}{\nobreak\hbox{-}}
\setlength{\emergencystretch}{2em}
\multlinegap0pt
\usepackage{bm}
\usepackage{float}
\usepackage[shortlabels]{enumitem}
\usepackage{xcolor}
\usepackage{amssymb}
\usepackage{mathtools}
\newtheorem{theorem}{Theorem}
\usepackage{cleveref}
\crefname{theorem}{Theorem}{Theorems}
\title{The Geometry of LLM Quantization: \\ GPTQ as Babai's Nearest Plane Algorithm}
\author{Jiale Chen, Yalda Shabanzadeh, Elvir Crn\v{c}evi\'c, Torsten Hoefler, Dan Alistarh}
\date{Source: arXiv:2507.18553v4 (CC BY 4.0)}
\begin{document}
\maketitle

\noindent\emph{Source and licence.} The Geometry of LLM Quantization: \\ GPTQ as Babai's Nearest Plane Algorithm. Version arXiv:2507.18553v4 (CC BY 4.0), distributed by arXiv under the Creative Commons Attribution 4.0 International licence, http://creativecommons.org/licenses/by/4.0/, as recorded in the arXiv licence metadata for this exact version and shown on the abstract page https://arxiv.org/abs/2507.18553v4. Full licence notice: \texttt{licenses/arxiv-2507.18553-CC-BY-4.0.txt}.\par\smallskip

\noindent\textit{Editorial scope.} Counted unit: the article's theorem thm:obq\_babai (source0.tex lines 426--430). The article's complete original proof of it is delivered with it (source0.tex lines 431--470, one \texttt{proof} environment of 40 lines). No mathematics is written, repaired or re-derived: the statement and the proof are the article's own lines, in the article's own notation, and no retained line is substituted or re-flowed.  Retained uncounted as the source-local context the proof needs: lines 188--191; lines 418--424.  One declared adaptation: exactly 1 ancillary strings inside the retained spans are omitted (image-inclusion commands and, where the source repeats a label, the repeated label definitions) and no other line differs from the source; the figure environments, with their placement, captions and labels, are kept, and the package records the omitted strings in fidelity receipts bound to the affected bodies.  No retained line contains a citation, so the wrapper carries no bibliography and no bibliography entry.  The retained lines contain the article's own diagram code, which the wrapper typesets with the standard tikz packages; no image file is delivered and the package declares no asset dependency.  Editorial formatting only, in addition: an AMS \texttt{amsart} preamble that restates the article's own declarations and macros used by the retained lines, namely the environments \texttt{theorem} on separate counters, together with the standard packages those lines need.  The retained environments print in this document as Theorem 1 in order of appearance, whereas the article prints the same items with the numbering of its own full document.  The complete native source of the article as distributed in the arXiv e-print for this exact version is bundled unchanged as \texttt{sources/182160/source0.tex}.
\begin{equation}
\Delta \bm{w}_{i} [j'] \leftarrow \frac{\left( \bm{X} [:, J]^\top \bm{X} [:, J] \right)^{-1}[j', j]}{\left( \bm{X} [:, J]^\top \bm{X} [:, J] \right)^{-1}[j, j]} \left( \bm{q}_{i} [j] - \bm{w}_{i} [j] \right)
\label{eq:obq_update} .
\end{equation}

\begin{figure}[ht]
\begin{center}

\end{center}
\caption{Equivalence of OBQ's error propagation and Babai's projection. \textbf{(a)} 3D plot showing the target being projected onto the nearest plane. \textbf{(b)} 3D plot showing how the projection error is propagated. \textbf{(c)} 2D plot showing the vectors on the nearest hyperplane in (a-b). \textbf{(d)} 2D plot showing the vectors on the orthogonal projection plane in (b).}
\label{fig:obq_3d}
\end{figure}

\begin{theorem}[Error Propagation and Babai's Projection]
\label{thm:obq_babai}
Babai's nearest plane algorithm iteratively projects the target vector onto the nearest hyperplane and rounds the coefficient.
The OBQ error propagation step (Eq.~\ref{eq:obq_update}) is exactly this projection on the original basis $\bm{B} = \bm{X} \mathrm{\ diag} \left( \bm{s}_{i} \right)$ without basis reduction.
\end{theorem}

\begin{proof}
Let $\bm{B} = \left[ \bm{b}_{1}, \dots, \bm{b}_{c} \right]$ be the basis with $\bm{b}_{j}$ being a basis vector.
Let $J$ be the set of unprojected indices with ${j}_{1}, {j}_{2} \in J$ and ${j}_{1} \ne {j}_{2}$.
Let $\bm{y} = \sum_{j \in J} \zeta_{j} \bm{b}_{j}$ be the current residual target where $\zeta_{j} \in \mathbb{R}$ is a real number to be rounded to integers.
Let $\mathcal{NHP} \coloneq \left\lfloor {\zeta}_{{j}_{2}} \right\rceil \bm{b}_{{j}_{2}} + \mathrm{Span} \left\{ \bm{b}_{j} \; | \; {j} \ne {j}_{2} \right\}$ be the nearest hyperplane that is orthogonal to the Gram-Schmidt vector $\bm{b}_{{j}_{2}} - \sum_{{j} \ne {j}_{2}} \mathrm{Proj}_{\bm{b}_{j}} \left( \bm{b}_{{j}_{2}} \right)$.
\cref{fig:obq_3d} (a) is a 3D plot showing the projection error vector $\Delta \bm{y} = \mathrm{Proj}_{\mathcal{NHP}} \left( \bm{y} \right) - \bm{y}$.
We focus on analyzing the error propagation in the direction of basis $\bm{b}_{{j}_{1}}$ induced by the projection of basis $\bm{b}_{{j}_{2}}$ and collapse the span of other basis vectors to a single dimension as illustrated by the hyperline $\mathcal{HL} \coloneq \left\lfloor {\zeta}_{{j}_{2}} \right\rceil \bm{b}_{{j}_{2}} + \mathrm{Span} \left\{ \bm{b}_{j} | {j} \ne {j}_{1}, {j}_{2} \right\}$.
\cref{fig:obq_3d} (b) is a 3D plot showing the decomposition of the error $\Delta \bm{y} = \sum_{j \in J} \Delta {\zeta}_{j} \bm{b}_{j}$ as the error component vectors in the basis directions.
\cref{fig:obq_3d} (c) is a 2D plot showing the vectors on plane $\mathcal{NHP}$.
The number ${\zeta}_{j}$ will be updated to ${\zeta}_{j} + \Delta {\zeta}_{j}$ such that $\mathrm{Proj}_{\mathcal{NHP}} \left( \bm{y} \right) = \sum_{j \in J} \left( \zeta_{j} + \Delta {\zeta}_{j} \right) \bm{b}_{j}$.
Next, let $\bm{N} = \bm{B}^{-\top} = \left[ \bm{n}_{1}, \dots, \bm{n}_{c} \right]$ be the inverse basis. Then, we have $\left< \bm{n}_{j} , \bm{b}_{j} \right> = 1$ and $\bm{n}_{j} \perp \bm{b}_{j'}, \forall j \ne j'$.
We project all the vectors in \cref{fig:obq_3d} (b) onto the orthogonal projection plane $\mathcal{OPP} \coloneq \mathrm{Span} \left\{ \bm{n}_{j} | j = {j}_{1}, {j}_{2} \right\}$ that is orthogonal to the hyperline $\mathcal{HL}$, and continue the proof in the 2D geometry in \cref{fig:obq_3d} (d).
Denote the angle $\theta = \angle \left( \bm{n}_{{j}_{1}}, \bm{n}_{{j}_{2}} \right) = \pi - \angle \left( \mathrm{Proj}_{\mathcal{OPP}} \left( \bm{b}_{{j}_{1}} \right), \mathrm{Proj}_{\mathcal{OPP}} \left( \bm{b}_{{j}_{2}} \right) \right)$.
Then, $
\frac{\Delta \zeta_{{j}_{1}} \left\| \mathrm{Proj}_{\mathcal{OPP}} \left( \bm{b}_{{j}_{1}} \right) \right\|}
{\Delta \zeta_{{j}_{2}} \left\| \mathrm{Proj}_{\mathcal{OPP}} \left( \bm{b}_{{j}_{2}} \right) \right\|}
= \cos \theta
= \frac{\left< \bm{n}_{{j}_{1}} , \bm{n}_{{j}_{2}} \right>}{\left\| \bm{n}_{{j}_{1}} \right\| \left\| \bm{n}_{{j}_{2}} \right\|}
= \frac{\left\| \bm{n}_{{j}_{2}} \right\|}{\left\| \bm{n}_{{j}_{1}} \right\|} \frac{\left< \bm{n}_{{j}_{1}} , \bm{n}_{{j}_{2}} \right>}{\left< \bm{n}_{{j}_{2}} , \bm{n}_{{j}_{2}} \right>}
$.
For ${j} = {j}_{1}, {j}_{2}$, $
\left\| \mathrm{Proj}_{\mathcal{OPP}} \left( \bm{b}_{j} \right) \right\| \left\| \bm{n}_{j} \right\|
= \frac{\left< \mathrm{Proj}_{\mathcal{OPP}} \left( \bm{b}_{j} \right) , \bm{n}_{j} \right>}{\cos \left( \frac{\pi}{2} - \theta \right)}
= \frac{\left< \bm{b}_{j} , \bm{n}_{j} \right>}{\cos \left( \frac{\pi}{2} - \theta \right)}
= \frac{1}{\cos \left( \frac{\pi}{2} - \theta \right)}
$.
For $j, j' \in \left\{ {j}_{1}, {j}_{2} \right\}$, $
\left< \bm{n}_{j} , \bm{n}_{j'} \right>
= \left( \bm{N}^{\top} \bm{N} \right) [j, j']
= \left( \bm{B}^{\top} \bm{B} \right)^{-1} [j, j']
$.
Combining the above equations,
$
\Delta \zeta_{{j}_{1}}
= \frac{\left\| \mathrm{Proj}_{\mathcal{OPP}} \left( \bm{b}_{{j}_{2}} \right) \right\| \left\| \bm{n}_{{j}_{2}} \right\|}{\left\| \mathrm{Proj}_{\mathcal{OPP}} \left( \bm{b}_{{j}_{1}} \right) \right\| \left\| \bm{n}_{{j}_{1}} \right\|} \frac{\left< \bm{n}_{{j}_{1}} , \bm{n}_{{j}_{2}} \right>}{\left< \bm{n}_{{j}_{2}} , \bm{n}_{{j}_{2}} \right>} \Delta \zeta_{{j}_{2}}
= \frac{\left< \bm{n}_{{j}_{1}} , \bm{n}_{{j}_{2}} \right>}{\left< \bm{n}_{{j}_{2}} , \bm{n}_{{j}_{2}} \right>} \Delta \zeta_{{j}_{2}}
= \frac{\left( \bm{B}^{\top} \bm{B} \right)^{-1} [{j}_{1}, {j}_{2}]}{\left( \bm{B}^{\top} \bm{B} \right)^{-1} [{j}_{2}, {j}_{2}]} \Delta \zeta_{{j}_{2}}
$.
Finally, substituting $\bm{B} = \left( \bm{X} \mathrm{\ diag} \left( \bm{s}_{i}\right) \right) [:, J]$ and ${\zeta}_{j} = \frac{\bm{w}_{i} [j]}{\bm{s}_{i} [j]}$ completes the proof.
\end{proof}

\end{document}
